% ECON 101: slide-faithful economics with elementary analysis explanations.
\documentclass[11pt,latinmodern]{article}
\usepackage[europe,paper=a4]{qhnotes}
% Tables and plots were already used in this note set.
\usepackage{booktabs}
\usepackage{pgfplots}
\pgfplotsset{compat=1.18}
\AtEndPreamble{%
  \RequirePackage[nameinlink]{cleveref}%
  \crefname{theorem}{Result}{Results}% shared qhnotes theorem counter
  \Crefname{theorem}{Result}{Results}%
  \crefname{example}{Example}{Examples}%
  \Crefname{example}{Example}{Examples}%
  \crefname{definition}{Definition}{Definitions}%
  \Crefname{definition}{Definition}{Definitions}%
}
\newcommand{\R}{\mathbb{R}}
\newcommand{\N}{\mathbb{N}}
\newcommand{\abs}[1]{\left|#1\right|}
\newcommand{\set}[1]{\left\{#1\right\}}
\newcommand{\paren}[1]{\left(#1\right)}
\newcommand{\dollar}[1]{\$#1}
\DeclareMathOperator{\TB}{TB}
\DeclareMathOperator{\TC}{TC}
\DeclareMathOperator{\MB}{MB}
\DeclareMathOperator{\MC}{MC}
\DeclareMathOperator{\FC}{FC}
\DeclareMathOperator{\VC}{VC}
\DeclareMathOperator{\MP}{MP}
\title{\textbf{ECON 101}\\[0.5em]\large Introduction to Microeconomics}
\author{Qinghao Hu\\\small Based on lectures by Mikal Skuterud}
\date{Fall 2026}
\course{Introduction to Microeconomics}
\begin{document}
\raggedbottom
% All statement boxes are short; keep headings and contents together.
\mdfsetup{nobreak=true}
\maketitle
\begin{quote}\small
These notes explain Chapters 1--3 of the lecture slides, preserving their
examples and numerical data. The aim is to make the reasoning and its
conditions explicit, rather than to introduce a separate mathematical
model. Definitions and qualifications are the note-taker's exposition;
the connecting lines in redrawn graphs are illustrative.
\end{quote}
\tableofcontents
%!TEX root = ../note.tex
% Source: ../chapter1_econ101_fall2026.pdf
\section{The four core principles of economics}
An \term{economy} is a system that answers four questions for a group of
agents facing scarce resources: \emph{what} to produce, \emph{how} to produce
it, \emph{who} does what, and \emph{who} gets what. (The slides' example is
the island in Golding's \emph{Lord of the Flies}.) The basic unit of analysis
is an \emph{individual decision}, and four principles tell us how decisions
are made:
\begin{enumerate}
  \item the cost--benefit principle,
  \item the opportunity cost principle,
  \item the marginal principle,
  \item the interdependence principle.
\end{enumerate}

\subsection{The cost--benefit principle}
Throughout, an agent chooses from a set $A$ of feasible options.

\begin{defi}
  For an option $a \in A$, let $B(a)$ be its \term{benefit} and $C(a)$ its
  \term{cost}, both measured in dollars. The \term{economic surplus} of $a$ is
  \[
    S(a) = B(a) - C(a).
  \]
\end{defi}

Benefits and costs need not be monetary. They are converted to dollars through
willingness to pay.

\begin{defi}
  The \term{willingness to pay} for a benefit is the largest amount of money
  $w$ the agent would give up to get it. For a cost, it is the largest $w$ the
  agent would pay to avoid it.
\end{defi}

\begin{principle}[Cost--benefit principle]
  Pursue an option $a$ if and only if $B(a) \geq C(a)$, i.e.\ $S(a) \geq 0$.
\end{principle}

\begin{remark}
  Some complications (the slides list them but do not resolve them yet):
  \begin{itemize}
    \item not everything is about money, which is why willingness to pay is
      used as the unit;
    \item the benefit may go to someone else;
    \item $B$ and $C$ may be uncertain, in which case they are random variables.
  \end{itemize}
  Costs and benefits are the \term{economic incentives} that drive decisions.
\end{remark}

\subsubsection*{Gains from trade}
\begin{nprop}[Voluntary exchange is not zero-sum]\label{prop:trade}
  A buyer values a good at $v$, a seller can produce it at cost $c$, and they
  trade at price $p$. Then
  \[
    S_{\text{buyer}} = v - p,\qquad S_{\text{seller}} = p - c,\qquad
    S_{\text{buyer}} + S_{\text{seller}} = v - c .
  \]
  Both agree to trade if and only if $c \leq p \leq v$, in which case both
  surpluses are nonnegative. The total surplus $v - c$ does not depend on $p$.
  The price only decides how it is split.
\end{nprop}
\begin{proof}
  The buyer accepts iff $v - p \geq 0$ and the seller accepts iff
  $p - c \geq 0$. Adding the two surpluses cancels $p$.
\end{proof}

\begin{significant}
  Voluntary exchange creates surplus $v - c \geq 0$. Nobody loses.
\end{significant}

\subsection{The opportunity cost principle}
\begin{defi}
  Given a choice $a \in A$, the \term{next best alternative} is the best
  option in $A \setminus \{a\}$. The \term{opportunity cost} of $a$ is the
  value of what is given up:
  \[
    \OC(a) = \max_{a' \in A \setminus \{a\}} (\text{value of } a').
  \]
  The out-of-pocket part of a cost is the \term{financial cost}. The
  opportunity cost can be much larger.
\end{defi}

\begin{principle}[Opportunity cost principle]
  The true cost of something is what you must give up to get it. When deciding
  on $a$, always ask ``$a$, \emph{or} the next best alternative?''
\end{principle}

Resources (income, time, energy) are \term{scarce}. That is why every choice
has an opportunity cost.

\begin{defi}
  A \term{sunk cost} is a cost that has to be paid whatever option is
  chosen, i.e.\ a cost $k$ that enters $C(a)$ for every $a \in A$.
\end{defi}

\begin{nprop}[Ignore sunk costs]\label{prop:sunk}
  If $C(a) = \tilde C(a) + k$ for all $a \in A$ with $k$ constant, then
  \[
    \argmax_{a \in A} \bigl(B(a) - C(a)\bigr) = \argmax_{a \in A} \bigl(B(a) - \tilde C(a)\bigr).
  \]
\end{nprop}
\begin{proof}
  Subtracting the constant $k$ from every surplus preserves their order.
\end{proof}

\begin{eg}
  Tuition is \dollar{4{,}500} per term for 5 courses with 24 lectures each, so
  one lecture costs $\frac{4500}{5 \cdot 24} = \dollar{37.50}$. You pay this
  whether you attend or not, so it is sunk and does not affect the decision.
  The relevant cost of going to lecture is the next best use of those 80
  minutes, e.g.\ wages from a job.
\end{eg}

\begin{eg}
  The biggest cost of an hour of gaming is not the hardware or software. It is
  the hour, which could have been spent on something else.
\end{eg}

\begin{ex}
  Explain each of the following with the opportunity cost principle.
  \begin{enumerate}
    \item Many Canadians date as teenagers but few marry as teenagers.
    \item There are fewer stay-at-home mothers than 60 years ago.
    \item More Canadians are born in September--October than in February--March.
    \item Terminally ill patients are more willing to try unproven drugs.
  \end{enumerate}
\end{ex}
\begin{remark}[Sketch]
  In (ii), higher wages for women raise the opportunity cost of staying at
  home. In (iv), the patient gives up fewer expected years of life, so the
  opportunity cost of the risk is lower.
\end{remark}

\subsubsection*{Production possibilities frontier}
\begin{defi}
  Given scarce resources, the \term{feasible set} $F \subseteq \R_{\geq 0}^2$ is
  the set of output pairs $(x, y)$ that can be produced. A point
  $(x, y) \in F$ is \term{efficient} if there is no $(x', y') \in F$ with
  $x' \geq x$, $y' \geq y$ and $(x', y') \neq (x, y)$. The
  \term{production possibilities frontier} (PPF) is the set of efficient
  points in $F$.
\end{defi}

On a PPF $y = f(x)$ with $f$ decreasing, the opportunity cost of one more
unit of $x$ is the amount of $y$ given up, $-\Delta y / \Delta x$. For a
differentiable $f$ this is $\abs{f'(x)}$, the absolute value of the slope.

\begin{eg}[Studying]\label{eg:ppf}
  You have $3$ hours a night for economics ($h_e$) and psychology ($h_p$), so
  $h_e + h_p \leq 3$. Each hour raises your economics grade by $8$ points and
  your psychology grade by $4$ points:
  \[
    E = 8h_e,\qquad P = 4h_p.
  \]
  Substituting $h_e = E/8$ and $h_p = P/4$ gives
  \[
    F = \set{(P, E) \in \R_{\geq 0}^2 : \tfrac{E}{8} + \tfrac{P}{4} \leq 3},
    \qquad \text{PPF: } E = 24 - 2P,\quad 0 \leq P \leq 12.
  \]
  The opportunity cost of one psychology point is $2$ economics points. For
  example, going from $(4, 16)$ to $(8, 8)$ costs $8$ economics points for
  $4$ psychology points. Points strictly inside $F$, such as $(3, 10)$, are
  \emph{inefficient}.
\end{eg}

\begin{center}
\begin{tikzpicture}
  \begin{axis}[width=7cm, height=6cm, xmin=0, xmax=20, ymin=0, ymax=38,
      xlabel={Psych boost $P$}, ylabel={Econ boost $E$},
      axis lines=left, xtick={4,8,12,18}, ytick={8,16,24,36}]
    \addplot[thick, red, domain=0:12] {24 - 2*x};
    \addplot[thick, blue, dashed, domain=0:18] {36 - 2*x};
    \addplot[only marks, mark=*] coordinates {(4,16) (8,8)};
    \addplot[only marks, mark=x, mark size=3pt] coordinates {(3,10)};
    \node[red] at (axis cs:3.5,4) {old PPF};
    \node[blue] at (axis cs:14,14) {new PPF};
  \end{axis}
\end{tikzpicture}
\end{center}

\begin{nprop}[Features of a PPF]
  \leavevmode
  \begin{enumerate}
    \item Every point on the PPF is efficient: one output can only be raised
      by lowering the other.
    \item The absolute value of the slope of the PPF is the opportunity cost
      of the good on the horizontal axis.
    \item Higher productivity enlarges $F$ and pushes the PPF outward.
  \end{enumerate}
\end{nprop}

\begin{eg}
  In Example~\ref{eg:ppf}, suppose productivity rises by $50\%$ in both
  subjects: $E = 12h_e$ and $P = 6h_p$. The new PPF is $E = 36 - 2P$ for
  $0 \leq P \leq 18$. The frontier moves out, but the opportunity cost is
  still $12/6 = 2$.
\end{eg}

\subsection{The marginal principle}
So far choices were binary. Many choices are about \emph{how much}: choose a
quantity $q \in \set{0, 1, 2, \ldots}$.

\begin{defi}
  Let $\TB(q)$ and $\TC(q)$ be the \term{total benefit} and \term{total cost}
  of quantity $q$. For $q \geq 1$ the \term{marginal benefit} and
  \term{marginal cost} of the $q$-th unit are
  \[
    \MB(q) = \TB(q) - \TB(q-1),\qquad \MC(q) = \TC(q) - \TC(q-1).
  \]
  The surplus is $S(q) = \TB(q) - \TC(q)$.
\end{defi}

\begin{principle}[Marginal principle]
  Break a decision about quantity into a sequence of marginal decisions:
  ``should I get one more?'' Answer each one with the cost--benefit
  principle, i.e.\ compare $\MB(q+1)$ with $\MC(q+1)$.
\end{principle}

\begin{nlemma}\label{lem:telescope}
  For all $q \geq 0$,
  \[
    S(q) = S(0) + \sum_{k=1}^{q} \bigl(\MB(k) - \MC(k)\bigr).
  \]
\end{nlemma}
\begin{proof}
  The sum telescopes:
  $\sum_{k=1}^q \MB(k) = \TB(q) - \TB(0)$, and similarly for $\MC$.
\end{proof}

\begin{nthm}[Rational rule]\label{thm:rational}
  Suppose $\MB - \MC$ is nonincreasing (this holds, for instance, if $\MB$
  is nonincreasing and $\MC$ is nondecreasing). Then
  \[
    q^* = \max\set{q \geq 0 : q = 0 \text{ or } \MB(q) \geq \MC(q)}
  \]
  maximises $S$. In words: \emph{keep going as long as the marginal benefit
  is at least the marginal cost.}
\end{nthm}
\begin{proof}
  Write $d(k) = \MB(k) - \MC(k)$, which is nonincreasing. By the choice of
  $q^*$ and monotonicity, $d(k) \geq 0$ for $k \leq q^*$ and $d(k) < 0$ for
  $k > q^*$. By Lemma~\ref{lem:telescope}, if $q < q^*$ then
  $S(q^*) - S(q) = \sum_{k=q+1}^{q^*} d(k) \geq 0$, and if $q > q^*$ then
  $S(q) - S(q^*) = \sum_{k=q^*+1}^{q} d(k) < 0$.
\end{proof}

\begin{remark}
  If $q$ can vary continuously and $\TB$, $\TC$ are differentiable, then
  $\MB = \TB'$ and $\MC = \TC'$. An interior optimum satisfies the
  first-order condition $S'(q^*) = 0$, i.e.
  \[
    \MB(q^*) = \MC(q^*).
  \]
  It is a maximum when $S'' = \MB' - \MC' \leq 0$. This is the
  ``marginal benefit equals marginal cost'' form of the rule.
\end{remark}

\begin{eg}[Nerida's restaurant]
  Each meal sells for \dollar{25}. Total cost is \dollar{10} per meal in food,
  \dollar{300} per waiter, \dollar{500} rent, and \dollar{1{,}000} for
  Nerida's own time. How many waiters should she hire?
  \begin{center}
    \begin{tabular}{ccccccc}
      \toprule
      Workers & Meals & $\TB$ & $\MB$ & $\TC$ & $\MC$ & $S$ \\
      \midrule
      2 & 160 & 4000 & ---  & 3700 & --- & 300 \\
      3 & 210 & 5250 & 1250 & 4500 & 800 & 750 \\
      4 & 250 & 6250 & 1000 & 5200 & 700 & 1050 \\
      5 & 280 & 7000 & 750  & 5800 & 600 & 1200 \\
      6 & 300 & 7500 & 500  & 6300 & 500 & \textbf{1200} \\
      7 & 310 & 7750 & 250  & 6700 & 400 & 1050 \\
      \bottomrule
    \end{tabular}
  \end{center}
  For example, the 4th worker adds $40$ meals, so
  $\MB(4) = 25 \cdot 40 = 1000$ and $\MC(4) = 300 + 10 \cdot 40 = 700$.
  Here both $\MB$ and $\MC$ fall, but $\MB - \MC = 450, 300, 150, 0, -150$
  is still decreasing, so Theorem~\ref{thm:rational} applies. The
  rule stops at $q^* = 6$, where $\MB = \MC = 500$, and the maximum surplus is
  \dollar{1{,}200}. (The 6th worker adds exactly zero surplus, so $q = 5$ is
  also optimal.)
\end{eg}

\begin{remark}
  Rent and Nerida's time are the same for every $q$. They are sunk with
  respect to this decision and never appear in $\MC$
  (Proposition~\ref{prop:sunk}).
\end{remark}

\subsection{The interdependence principle}
\begin{principle}[Interdependence principle]
  Your best choice in one situation depends on:
  \begin{enumerate}
    \item \emph{your other choices}: resources (income, time, attention,
      wealth) are shared across all your decisions;
    \item \emph{other people's choices}: scarcity creates competition between
      agents (people, businesses, governments);
    \item \emph{other markets}: e.g.\ labour, housing and credit markets affect
      one another;
    \item \emph{time}: today's choice depends on past choices and expected
      future ones.
  \end{enumerate}
\end{principle}

\begin{remark}
  Mathematically, (i) says you solve one problem subject to a joint budget
  constraint (as in the PPF), not several problems separately. Points (ii) and
  (iii) say that the parameters of your problem, such as prices, are
  determined by the choices of everyone else.
\end{remark}

%!TEX root = ../note.tex
% Source: ../chapter2_econ101_fall2026.pdf
\section{Demand and consumer choice}
\subsection{Individual demand}
\begin{defi}
  The \term{individual demand curve} of a buyer $i$ is the function
  \[
    q_i^D : \R_{\geq 0} \to \R_{\geq 0},\qquad p \mapsto q_i^D(p),
  \]
  giving the quantity demanded at each price $p$, \emph{holding everything
  else constant} (\term{ceteris paribus}). By convention price goes on the
  vertical axis and quantity on the horizontal axis.
\end{defi}

\begin{eg}[Darren's demand for gasoline]\label{eg:darren}
  A survey of Darren's daily gas use (litres) at hypothetical prices:
  \begin{center}
    \begin{tabular}{cccccc}
      \toprule
      $p$ (\$/litre) & 2.40 & 2.20 & 2.00 & 1.80 & 1.60 \\
      $q^D_{\text{Darren}}(p)$ & 1 & 2 & 3 & 5 & 7 \\
      \bottomrule
    \end{tabular}
  \end{center}
\end{eg}

\begin{law}[Law of demand]
  Quantity demanded rises when the price falls, and vice versa:
  \[
    p_1 < p_2 \implies q^D(p_1) \geq q^D(p_2).
  \]
  That is, demand curves are \emph{downward sloping}.
\end{law}

\subsection{The rational rule for buyers}
To explain the law of demand, apply the marginal principle to a buyer who
takes the price $p$ as given. Buying the $k$-th unit has marginal benefit
$\MB(k)$ and marginal cost $p$.

\begin{nthm}[Rational rule for buyers]\label{thm:buyer}
  Suppose $\MB(1) \geq \MB(2) \geq \cdots$ (\term{diminishing marginal
  benefit}). At price $p$ the surplus-maximising quantity is
  \[
    q^D(p) = \max\set{k \geq 0 : k = 0 \text{ or } \MB(k) \geq p}
    = \#\set{k \geq 1 : \MB(k) \geq p}.
  \]
  At an interior optimum, $\text{price} = \text{marginal benefit}$.
\end{nthm}
\begin{proof}
  Apply Theorem~\ref{thm:rational} with $\MC(k) = p$ for all $k$. Since $\MB$
  is nonincreasing, the set $\set{k : \MB(k) \geq p}$ is an initial segment
  $\set{1, \ldots, q^D(p)}$, so its largest element equals its size.
\end{proof}

\begin{ncor}\label{cor:lawdemand}
  Diminishing marginal benefit implies the law of demand. Moreover, the
  demand curve \emph{is} the marginal benefit curve.
\end{ncor}
\begin{proof}
  If $p_1 < p_2$ then $\set{k : \MB(k) \geq p_2} \subseteq \set{k : \MB(k)
  \geq p_1}$, so $q^D(p_2) \leq q^D(p_1)$ by Theorem~\ref{thm:buyer}.
\end{proof}

\begin{eg}\label{eg:darren-mb}
  Darren ranks the uses of each extra litre:
  \begin{center}
    \begin{tabular}{clc}
      \toprule
      $k$ & Use of the $k$-th litre & $\MB(k)$ \\
      \midrule
      1 & drive to work instead of the bus & 2.40 \\
      2 & drive to school instead of biking & 2.20 \\
      3 & shop at the cheaper No Frills & 2.00 \\
      4 & drive to see friends & 1.90 \\
      5 & drive to the gym & 1.80 \\
      6 & drive to do errands & 1.70 \\
      7 & scenic drive & 1.60 \\
      \bottomrule
    \end{tabular}
  \end{center}
  At $p = 1.85$, $\set{k : \MB(k) \geq 1.85} = \set{1, 2, 3, 4}$, so Darren
  buys $4$ litres a day. The formula also reproduces the survey in
  Example~\ref{eg:darren}: e.g.\ at $p = 1.80$ he buys $5$ litres.
\end{eg}

\begin{center}
\begin{tikzpicture}
  \begin{axis}[width=7cm, height=5.5cm, xmin=0, xmax=8, ymin=1.5, ymax=2.5,
      xlabel={Quantity (litres/day)}, ylabel={Price / $\MB$ (\$)},
      axis lines=left, xtick={1,...,7}, ytick={1.6,1.8,2.0,2.2,2.4},
      yticklabel style={/pgf/number format/fixed, /pgf/number format/precision=2}]
    \addplot[const plot, thick, blue] coordinates
      {(0,2.4) (1,2.2) (2,2.0) (3,1.9) (4,1.8) (5,1.7) (6,1.6) (7,1.6)};
    \addplot[dashed] coordinates {(0,1.85) (8,1.85)};
    \node[anchor=west] at (axis cs:4.2,1.9) {$p = 1.85$};
  \end{axis}
\end{tikzpicture}
\end{center}

\subsection{Market demand}
\begin{defi}
  If the buyers are $i = 1, \ldots, N$, the \term{market demand curve} is the
  horizontal sum of the individual demand curves:
  \[
    Q^D(p) = \sum_{i=1}^{N} q_i^D(p).
  \]
\end{defi}

\begin{nprop}
  If every $q_i^D$ is nonincreasing, then so is $Q^D$.
\end{nprop}
\begin{proof}
  A sum of nonincreasing functions is nonincreasing.
\end{proof}

A movement along $Q^D$ reflects two things: existing buyers changing the
quantity they buy, and buyers entering or leaving the market ($q_i^D(p)$
changing between $0$ and a positive value).

\begin{eg}[Estimating market demand]
  Survey a random sample of $n = 400$ people from a population of
  $N = 40$ million. If $\bar q(p)$ is the average demand in the sample, then
  $Q^D(p) \approx N \bar q(p) = \frac{N}{n} \sum_{i \in \text{sample}}
  q_i^D(p)$, with scale factor $\frac{N}{n} = 100{,}000$.
  \begin{center}
    \begin{tabular}{cccc}
      \toprule
      $p$ & Sample total (litres) & $\times\, N/n$ & $Q^D(p)$ (millions) \\
      \midrule
      1.60 & 1{,}200 & $10^5$ & 120 \\
      1.80 & 1{,}100 & $10^5$ & 110 \\
      2.00 & 1{,}000 & $10^5$ & 100 \\
      2.20 & 900     & $10^5$ & 90 \\
      2.40 & 800     & $10^5$ & 80 \\
      \bottomrule
    \end{tabular}
  \end{center}
\end{eg}

\subsection{Movements along versus shifts of demand}
Demand really depends on many variables:
\[
  Q^D = Q^D(p;\ I,\ \theta,\ p_r,\ e,\ \ldots,\ N),
\]
where $I$ is income, $\theta$ preferences, $p_r$ prices of related goods,
$e$ expectations and $N$ the number of buyers. The demand \emph{curve} is
the graph of $p \mapsto Q^D(p; \cdot)$ with all other arguments fixed.

\begin{defi}
  A change in the good's own price $p$ is a \term{movement along the demand
  curve} and changes the \term{quantity demanded}. A change in any other
  argument is a \term{shift of the demand curve} and changes
  \term{demand}. An \emph{increase in demand} is a shift to the right (more
  demanded at every $p$). A \emph{decrease} is a shift to the left.
\end{defi}

\begin{warning}
  ``The price fell, so demand rose'' is wrong. A price fall raises the
  \emph{quantity demanded}; the curve itself does not move. For example, a
  price drop from \dollar{2.20} to \dollar{1.80} moves the market along the
  curve from $90$ to $110$ million litres.
\end{warning}

\subsubsection*{The six demand shifters}
\begin{enumerate}
  \item \emph{Income.}
    \begin{defi}
      A good is \term{normal} if demand rises with income
      ($\partial Q^D / \partial I > 0$) and \term{inferior} if it falls
      ($\partial Q^D / \partial I < 0$).
    \end{defi}
  \item \emph{Preferences}, e.g.\ a marketing campaign convincing you that a
    brand is ``cool''.
  \item \emph{Prices of related goods.}
    \begin{defi}
      Let $p_y$ be the price of another good $y$. Then $y$ is a
      \term{complement} of $x$ if $\partial Q_x^D / \partial p_y < 0$ (a
      cheaper $y$ raises demand for $x$), and a \term{substitute} if
      $\partial Q_x^D / \partial p_y > 0$ (a cheaper $y$ lowers demand for
      $x$).
    \end{defi}
  \item \emph{Expectations}: beliefs about future prices affect demand today.
  \item \emph{Congestion and network effects.}
    \begin{defi}
      There is a \term{network effect} if a good is worth more to each buyer
      when more other people use it, and a \term{congestion effect} if it is
      worth less.
    \end{defi}
  \item \emph{Type and number of buyers.} Since $Q^D = \sum_{i=1}^N q_i^D$,
    raising $N$ shifts the \emph{market} curve right. Unlike factors (i)--(v),
    it does not shift any individual curve $q_i^D$.
\end{enumerate}

\begin{eg}
  A cut in income tax rates raises consumers' disposable income $I$. If gas
  is a normal good, market demand for gas shifts to the right: at every
  price, more gas is demanded.
\end{eg}

\begin{ex}
  Classify each change as a movement along or a shift of the demand for
  gasoline, and give the direction: (a) the price of gas falls; (b) the price
  of electric cars (a substitute for gas-powered driving) falls; (c) the
  population grows.
\end{ex}
\begin{remark}[Answers]
  (a) Movement down along the curve, so quantity demanded rises. (b) Left
  shift, since the price of a substitute fell. (c) Right shift of the market
  curve (factor (vi)).
\end{remark}

%!TEX root = ../note.tex
% Source: ../chapter3_econ101_fall2026.pdf
\section{Supply and producer choice}
This chapter mirrors Chapter 2 with sellers in place of buyers.

\subsection{Individual supply}
\begin{defi}
  The \term{individual supply curve} of a seller $j$ is the function
  $p \mapsto q_j^S(p)$ giving the quantity supplied at each price $p$,
  ceteris paribus.
\end{defi}

\begin{eg}[Shell]\label{eg:shell}
  Shell's supply of gas (millions of litres per day):
  \[
    q^S_{\text{Shell}}(p) =
    \begin{cases}
      0 & p < 1.60,\\
      8 + 5(p - 1.60) & 1.60 \leq p \leq 2.40,
    \end{cases}
  \]
  so $q^S(1.60) = 8$, $q^S(1.80) = 9$, $q^S(2.00) = 10$, $q^S(2.20) = 11$ and
  $q^S(2.40) = 12$. Below \dollar{1.60} Shell stops supplying gas.
\end{eg}

\begin{law}[Law of supply]
  Quantity supplied rises when the price rises, and vice versa:
  \[
    p_1 < p_2 \implies q^S(p_1) \leq q^S(p_2).
  \]
  That is, supply curves are \emph{upward sloping}.
\end{law}

\subsection{Perfect competition}
\begin{defi}
  A market is \term{perfectly competitive} if
  \begin{enumerate}
    \item there are many buyers and sellers, each small relative to the
      market, and no barriers to entry or exit;
    \item all firms sell an identical product;
    \item there is perfect information;
    \item mobility costs are zero: buyers can switch sellers for free.
  \end{enumerate}
  An agent is a \term{price taker} if it treats the market price $p$ as
  given, as opposed to a \emph{price maker}.
\end{defi}

\begin{nprop}
  In a perfectly competitive market every firm sells at the market price $p$.
\end{nprop}
\begin{proof}[Argument]
  If a firm charges $p' > p$, then by (ii)--(iv) every buyer switches to a
  rival at no cost, so it sells nothing. Suppose instead that selling at some
  $p' < p$ were profitable. With perfect information and free entry, a
  competitor would already have taken that opportunity, so this cannot hold in
  a situation where nobody wants to change what they do. Hence $p$ is
  exogenous to each firm.
\end{proof}

\begin{remark}
  Perfectly competitive markets are rare in reality. The assumption is a
  simplification that makes the analysis tractable.
\end{remark}

\subsection{The rational rule for sellers}
\begin{defi}
  A firm's total cost splits as
  \[
    \TC(q) = \FC + \VC(q),
  \]
  where the \term{fixed cost} $\FC$ does not depend on output (e.g.\
  capital) and the \term{variable cost} $\VC(q)$ does (e.g.\ labour). The
  \term{marginal cost} is $\MC(q) = \TC(q) - \TC(q-1)$, or $\TC'(q)$ in the
  continuous case.
\end{defi}

\begin{nprop}
  $\MC(q) = \VC(q) - \VC(q-1)$. In particular, fixed costs do not affect the
  choice of quantity.
\end{nprop}
\begin{proof}
  $\FC$ cancels in the difference. This is Proposition~\ref{prop:sunk} again.
\end{proof}

For a price taker, the marginal benefit of selling one more unit is the price
$p$.

\begin{nthm}[Rational rule for sellers]\label{thm:seller}
  If $\MC$ is nondecreasing, the profit-maximising quantity at price $p$ is
  \[
    q^S(p) = \max\set{q \geq 0 : q = 0 \text{ or } \MC(q) \leq p}.
  \]
  At an interior optimum, $\text{price} = \text{marginal cost}$. So:
  \emph{sell one more unit if $p \geq \MC$.}
\end{nthm}
\begin{proof}
  Apply Theorem~\ref{thm:rational} with $\MB(q) = p$ constant.
\end{proof}

\begin{ncor}
  Increasing marginal cost implies the law of supply, and the supply curve
  \emph{is} the marginal cost curve.
\end{ncor}
\begin{proof}
  If $p_1 < p_2$ then $\set{q : \MC(q) \leq p_1} \subseteq \set{q : \MC(q)
  \leq p_2}$.
\end{proof}

\begin{remark}
  The law of supply comes from two sources: \emph{diminishing marginal
  product} (below) and \emph{rising input costs}.
\end{remark}

\subsection{Diminishing marginal product}
\begin{defi}
  A \term{production function} $Y = F(L, K)$ gives output from labour $L$
  (variable input) and capital $K$ (fixed input). The \term{marginal
  product of labour} is $\MP_L = \partial F / \partial L$, or
  $F(L, K) - F(L-1, K)$ in discrete terms. There is \term{diminishing
  marginal product} if $\MP_L$ decreases in $L$ for fixed $K$.
\end{defi}

\begin{nthm}[Diminishing marginal product $\Rightarrow$ rising marginal cost]
  \label{thm:mpmc}
  Let the wage be $w$ and $K$ fixed, and let $F(\cdot, K)$ be smooth and
  strictly increasing, with inverse $L(Y)$ (the labour needed to make $Y$).
  Then
  \[
    \MC(Y) = \frac{w}{\MP_L\bigl(L(Y)\bigr)}.
  \]
  Hence if $\MP_L$ is decreasing, $\MC$ is increasing.
\end{nthm}
\begin{proof}
  $\TC(Y) = wL(Y) + \FC$, so by the inverse function rule
  $\MC(Y) = wL'(Y) = w / F_L(L(Y), K)$. When $Y$ increases, $L(Y)$ increases,
  so $\MP_L$ falls and $\MC$ rises.
\end{proof}

\begin{significant}
  Each extra unit of output needs more and more labour, so each unit costs
  more than the last.
\end{significant}

\begin{eg}
  Take $Y = \sqrt{LK}$ with $K = 4$, so $Y = 2\sqrt{L}$ and
  $\MP_L = \frac12\sqrt{K/L} = 1/\sqrt{L}$, which is decreasing. In discrete
  terms:
  \begin{center}
    \begin{tabular}{cccccccc}
      \toprule
      $L$ & 0 & 1 & 2 & 3 & 4 & 5 & 6 \\
      $Y$ & 0 & 2 & 2.83 & 3.46 & 4.00 & 4.47 & 4.90 \\
      $\MP_L$ & --- & 2 & 0.83 & 0.64 & 0.54 & 0.47 & 0.43 \\
      \bottomrule
    \end{tabular}
  \end{center}
  Invert: $L = Y^2 / K = Y^2/4$. With $w = \dollar{15}$,
  $\TC(Y) = 15 \cdot Y^2/4 = 3.75\,Y^2$ (ignoring the fixed cost of capital),
  so
  \[
    \MC(Y) = \TC(Y) - \TC(Y-1) = 3.75\,(2Y - 1).
  \]
  \begin{center}
    \begin{tabular}{cccccccc}
      \toprule
      $Y$ & 0 & 1 & 2 & 3 & 4 & 5 & 6 \\
      $L$ & 0 & 0.25 & 1 & 2.25 & 4 & 6.25 & 9 \\
      $\TC$ & 0 & 3.75 & 15 & 33.75 & 60 & 93.75 & 135 \\
      $\MC$ & --- & 3.75 & 11.25 & 18.75 & 26.25 & 33.75 & 41.25 \\
      \bottomrule
    \end{tabular}
  \end{center}
  In the continuous version, $\MC(Y) = \TC'(Y) = 7.5\,Y$, which matches
  Theorem~\ref{thm:mpmc}: $w / \MP_L(L(Y)) = 15\sqrt{Y^2/4} = 7.5\,Y$. Setting
  $p = \MC$ gives the supply curve $q^S(p) = p / 7.5$, which is upward
  sloping.
\end{eg}
\begin{remark}
  The table's $L$ column assumes labour can be hired in fractions. $\MP_L$
  and $\MC$ move in opposite directions, as Theorem~\ref{thm:mpmc} predicts.
\end{remark}

\subsection{Market supply}
\begin{defi}
  The \term{market supply curve} is the horizontal sum
  $Q^S(p) = \sum_{j=1}^{M} q_j^S(p)$ over the $M$ sellers. If the sellers
  are identical, $Q^S(p) = M\,q^S(p)$.
\end{defi}

\begin{eg}
  With $M = 10$ identical firms:
  \begin{center}
    \begin{tabular}{cccccc}
      \toprule
      $p$ & 1.60 & 1.80 & 2.00 & 2.20 & 2.40 \\
      $q^S(p)$ & 6 & 8 & 10 & 12 & 14 \\
      $Q^S(p) = 10\,q^S(p)$ & 60 & 80 & 100 & 120 & 140 \\
      \bottomrule
    \end{tabular}
  \end{center}
  (millions of litres per day).
\end{eg}

A higher price raises market quantity supplied in two ways: existing firms
expand output, and new firms enter.

\subsection{Movements along versus shifts of supply}
As with demand, write
\[
  Q^S = Q^S(p;\ w,\ A,\ p_{\text{rel}},\ e,\ M),
\]
where $w$ is input prices, $A$ productivity, $p_{\text{rel}}$ prices of
related outputs, $e$ expectations and $M$ the number of sellers. A change in
$p$ is a \term{movement along the supply curve} and changes the
\term{quantity supplied}. A change in any other argument is a \term{shift
of the supply curve}. An increase in supply is a shift to the right.

\subsubsection*{The five supply shifters}
\begin{enumerate}
  \item \emph{Input prices.} Higher input prices raise $\MC$ at every $q$, so
    supply shifts left.
  \item \emph{Productivity and technology.} Producing more with the same
    inputs (new machinery, better management or processes, learning by doing)
    lowers $\MC$, so supply shifts right.
  \item \emph{Prices of related outputs.}
    \begin{defi}
      Good $y$ is a \term{substitute-in-production} for $x$ if it is made
      with the same inputs, and a \term{complement-in-production} if it is
      made together with $x$. Then
      \[
        \frac{\partial Q^S_x}{\partial p_y} < 0 \text{ for substitutes},\qquad
        \frac{\partial Q^S_x}{\partial p_y} > 0 \text{ for complements}.
      \]
    \end{defi}
  \item \emph{Expectations.} If firms expect the price to rise, they store
    output rather than sell it now, so supply today shifts left. This matters
    more for non-perishable goods (gasoline) than perishable ones (fresh
    fish).
  \item \emph{Type and number of sellers.} Changing $M$ shifts the market
    curve but no individual curve. For example, a Competition Bureau ruling on
    a merger changes $M$.
\end{enumerate}

\begin{eg}
  Crude oil is an input to gasoline. If its price rises, $\MC$ rises at every
  output level, so the market supply of gasoline shifts \emph{left}: at every
  price, less is supplied. On the slides, supply at \dollar{2.00} falls from
  $100$ to $60$ million litres.
\end{eg}

\subsection*{Summary: buyers versus sellers}
\begin{center}
  \begin{tabular}{lll}
    \toprule
    & Buyer & Seller \\
    \midrule
    Marginal benefit & $\MB(q)$ & $p$ \\
    Marginal cost & $p$ & $\MC(q)$ \\
    Rational rule & $p = \MB$ & $p = \MC$ \\
    Law & demand slopes down & supply slopes up \\
    Reason & diminishing $\MB$ & diminishing $\MP_L$, rising input costs \\
    Market curve & $\sum_i q_i^D$ & $\sum_j q_j^S$ \\
    \bottomrule
  \end{tabular}
\end{center}

\printindex
\end{document}
